%%%%%%%%%%%%%%%%%%%%%%%%%%%%%%%%%%%%%%%%%%%%%%%%%%%%%%%
%% Bachelor's & Master's Thesis Template             %%
%% Copyleft by Artur M. Brodzki & Piotr Woźniak      %%
%% Covered for Lab/Assigment report by Tymon Żarski  %%
%% Faculty of Electronics and Information Technology %%
%% Warsaw University of Technology, 2019-2020        %%
%%%%%%%%%%%%%%%%%%%%%%%%%%%%%%%%%%%%%%%%%%%%%%%%%%%%%%%

\documentclass[
    left=2.5cm,         % Sadly, generic margin parameter
    right=2.5cm,        % doesnt't work, as it is
    top=2.5cm,          % superseded by more specific
    bottom=3cm,         % left...bottom parameters.
    bindingoffset=6mm,  % Optional binding offset.
    nohyphenation=false % You may turn off hyphenation, if don't like.
]{eiti/eiti-report}

\langpol % Dla języka angielskiego mamy \langeng
\graphicspath{{img/}}             % Katalog z obrazkami.
\addbibresource{bibliografia.bib} % Plik .bib z bibliografią

\begin{document}

%--------------------------------------
% Strona tytułowa
%--------------------------------------
\instytut{Telekomunikacji i Cyberbezpieczeństwa}
\przedmiot{PPMGR - Pracownia Problemowa Magisterska}
%\specjalnosc{XXXXXX}
\title{
    Badanie bezpieczeństwa agentów AI w aplikacjach
}
\reportDescription{Analiza literatury i sytuacji na rynku protokołu MCP}

\author{Andrzej Pultyn}
\album{325213}
\prowadzacy{dr inż. Sławomir Kukliński}
\date{\today}
\maketitle

%--------------------------------------
% Spis treści
%--------------------------------------
\tableofcontents

%--------------------------------------
% Rozdziały
%--------------------------------------
\cleardoublepage % Zaczynamy od nieparzystej strony
\pagestyle{headings}

\newpage % Rozdziały zaczynamy od nowej strony

\section{Wprowadzenie}

Protokół MCP (ang. \textit{Model Context Protocol}) jest standardem \textit{open-source}, udostępnionym przez firmę Anthropic w listopadzie 2024 roku \cite{anthropic-intro}. Główną motywacją jego powstania było uproszczenie integracji dużych modeli językowych (LLM - ang. \textit{Large Language Models}) z zewnętrznymi źródłami danych. Dzisiaj protokół wykorzystywany jest do rozszerzania możliwości modeli o dodatkowe narzędzia. Dzięki nim dzisiejsi agenci AI potrafią o wiele więcej niż odpowiadanie na pytania, bazując na wiedzy zdobytej w wyniku trenowania modelu.

%--------------------------------------------
% Literatura
%--------------------------------------------
\cleardoublepage % Zaczynamy od nieparzystej strony
\printbibliography



%--------------------------------------------
% Spisy (opcjonalne)
%--------------------------------------------
\pagestyle{plain}

\vspace{0.8cm}

\vspace{1cm}          % vertical space
\listofappendicestoc  % Spis załączników

% Załączniki
\newpage
\appendix{Nazwa załącznika 1}
\appendix{Nazwa załącznika 2}
\appendix{Nazwa załącznika 3}



\end{document}
